% ==============================================================
\section{おわりに}\label{sec:conclusion}
% ==============================================================

本論文では，密集区間マイニングフレームワークの二つの対称的拡張として
\emph{反密集区間}と\emph{コントラスト密集パターン}を提案した．

反密集区間は密集区間の双対概念を形式化し，
パターンの持続的な\emph{不在}期間を捉える．
効率的なトップダウン探索戦略を可能にする逆単調性を証明し，
反密集区間が密集区間と同じ計算量で算出できることを示した．

コントラスト密集パターンは，時間レジーム間で密集区間構成を比較するための
構造化されたフレームワークを提供する．
四種類のトポロジー変化（出現・消失・増幅・収縮）の分類体系は，
パターンが時間とともにどのように変遷するかの解釈可能な特徴づけを与える．
付随する置換検定はBenjamini--Hochberg補正と組み合わせることで
統計的厳密性を確保する．

合成データを用いた実験評価により，反密集区間の効果的な復元（F1 = 0.789）と
信頼性の高いトポロジー変化分類（76\%の正解率）を実証した．
シミュレートされたキャンペーン終了シナリオへの適用では，
キャンペーン依存パターンを高い統計的有意性（$p = 0.005$）で特定できること，
また非影響パターンを安定と正しく分類できることを示した．

\paragraph{制約}
本研究には以下の制約がある．
第一に，現在のフレームワークはレジーム境界$\tau$の事前指定を必要とする．
実応用ではレジーム境界が自明でない場合が多く，
この仮定は適用範囲を制限する．
第二に，置換検定はタイムスタンプの交換可能性を仮定しており，
強い時間的自己相関が存在する場合にはこの仮定が破られ，
$p$値が過小推定される可能性がある．
第三に，実験評価は合成データに限定されており，
実世界データセットでの有効性は未検証である．
これはトップカンファレンス投稿に向けて最も優先的に対処すべき制約である．
第四に，増幅・収縮の分類精度（E2で76\%）は改善の余地があり，
マージン$\delta$の選択に関するガイドラインが未整備である．
第五に，反密集区間の概念はアルゴリズム的には密集区間の閾値反転であり，
概念的新規性は主にコントラスト密集パターンのトポロジー分類体系と
統計検定フレームワークにある．

\paragraph{今後の課題}
以下の具体的な方向性がさらなる研究に値する：

\emph{(i) レジーム境界の自動発見}：
現在のフレームワークは境界$\tau$の事前指定を要求する．
PELT~\cite{killick2012pelt}やADWIN~\cite{bifet2007adwin}等の
自動変化点検出との統合により，サポート時系列から
レジーム境界を自動的に推定し，全パターンに一括適用する二段階手法が考えられる．
具体的には，全パターンのサポート時系列の主成分に対する変化点検出を
第一段階とし，検出された変化点をコントラスト分析の境界として使用する．

\emph{(ii) 多レジーム比較}：
二レジーム比較を$k > 2$レジームに拡張し，
$\binom{k}{2}$個のペアワイズ比較に対する
Benjamini--Hochberg補正を適用する．
これにより，四半期ごとの季節比較（$k = 4$）や
月次トレンド分析（$k = 12$）が実現可能となる．

\emph{(iii) 実世界データセットでの検証}：
Dunnhumbyデータセット（約250万トランザクション，89週間の小売データ）における
キャンペーン前後のコントラスト分析，および
UCI Online Retailデータセットにおける季節性パターンの変化検出を
最優先の検証課題とする．
医療分野では，MIMICデータベースの処方パターンに対する
ガイドライン改定前後のコントラスト分析が有望な応用である．

\emph{(iv) 効率的な枝刈り戦略}：
反密集区間の逆単調性（定理~\ref{thm:anti_monotone}）を活用した
トップダウン枝刈りの実装と，ボトムアップ密集区間検出との
統合戦略の計算量評価．

\emph{(v) コントラスト結果の冗長性制御}：
閉パターン・極大パターンへの拡張により，
報告されるコントラスト密集パターン数を制御する．
密集区間の包含関係を利用した冗長性除去フィルタにより，
解釈可能性を維持しつつ出力サイズを削減する．

\emph{(vi) パターンライフサイクル分析}：
密集から反密集への遷移ダイナミクスを時系列として可視化し，
パターンの「誕生」「成長」「衰退」「消滅」のライフサイクルを
コントラスト分類の時間的系列として追跡する枠組みの構築．
